\documentclass[11pt,a4paper]{article}
\usepackage[margin=25mm]{geometry}
\usepackage{fontspec}
\setmainfont{TeX Gyre Pagella}
\usepackage{amsmath,amssymb,mathtools}
\usepackage{graphicx}
\usepackage{xcolor}
\usepackage[unicode,colorlinks=true,linkcolor=blue,urlcolor=blue]{hyperref}
\usepackage{bookmark}
\usepackage{enumitem}
\setlist{itemsep=3pt,topsep=5pt}
\setlength{\emergencystretch}{3em}
\newcommand{\tightlist}{\setlength{\itemsep}{0pt}\setlength{\parskip}{0pt}}
\newcommand{\passthrough}[1]{#1}
\newcommand{\pandocbounded}[1]{#1}
\makeatletter
\def\maxwidth{\ifdim\Gin@nat@width>\linewidth\linewidth\else\Gin@nat@width\fi}
\def\maxheight{\ifdim\Gin@nat@height>0.75\textheight0.75\textheight\else\Gin@nat@height\fi}
\makeatother
\setkeys{Gin}{width=\maxwidth,height=\maxheight,keepaspectratio}
\hypersetup{pdfauthor={OpenAI GPT-6.1 Sol, Codex, Ultra effort}}
\begin{document}
\section{Frobenius lifts and abstract
reciprocity}\label{frobenius-lifts-and-abstract-reciprocity}

\emph{Written by OpenAI GPT-6.1 Sol in Codex, Ultra effort, October
2026. Self-checked by the writing AI; no independent review is claimed.
Public domain (CC0).}

The quotient by norms is the arithmetic object we want to understand.
For an unramified extension of degree \(n\), valuation already labels
its classes by \(\mathbf Z/n\mathbf Z\), and Frobenius supplies the
matching cyclic generator. A ramified extension asks for something more:
how can an automorphism still produce a class when it is not itself a
residue-field Frobenius?

We answer that question by making the automorphism into Frobenius over a
different finite field inside an enlarged unramified tower. The proposed
class is then the norm of a prime element. The proof has three distinct
jobs: calibrate that prescription on unramified extensions, show that
choices and products change it only by norms, and establish the two
different functorial operations. These jobs determine the order below.

Only the closed-subgroup correspondence from
\href{../profinite-groups-and-infinite-galois-theory.html}{Profinite
groups and infinite Galois theory} and the cyclic Tate calculations from
\href{../cohomology-of-cyclic-groups-and-the-herbrand-quotient.html}{Cohomology
of cyclic groups and the Herbrand quotient} are required for the
abstract proof. The local arithmetic module is verified in
\href{../local-reciprocity-and-norm-groups.html}{Local reciprocity and
norm groups}, sections 1--3; those computations may be read first for
motivation. Sections 4--5 of that lesson apply the abstract theorem
after it has been proved. This separates the example from the proof it
is meant to illustrate.

\subsection{1. What a reciprocity class must
do}\label{what-a-reciprocity-class-must-do}

Consider a cyclic unramified local extension \(L/K\) of degree \(n\).
Unit norms are surjective and \(v_K(Nb)=n v_L(b)\). Therefore the
valuation map gives \[
K^\times/NL^\times\simeq\mathbf Z/n\mathbf Z,
\qquad [\pi_K]\longmapsto1.
\] For the kernel assertion, if \(v_K(a)=nm\), divide \(a\) by the norm
of \(\pi_K^m\), viewed in \(L\). The remainder is a unit, hence a unit
norm. This calculation uses only the unramified norm theorem, not
reciprocity. It suggests the unique map taking arithmetic Frobenius to
\([\pi_K]\).

There are two different changes of field to track. Restricting an
automorphism should correspond to taking a norm of the arithmetic
element. Including an arithmetic element in a larger ground field should
correspond to transfer of the automorphism. For example, in an
unramified cyclic extension of degree six, the degree-three intermediate
field has Frobenius equal to the third power of ground-field Frobenius.
Transfer from the order-six group to its order-two subgroup sends a
generator to that third power. On multiplicative groups a ground-field
uniformizer stays a uniformizer, whereas its norm down a degree-three
unramified extension has valuation three. Thus the two operations cannot
be interchanged.

We construct the map from automorphisms to norm classes. The next lesson
proves that it is an isomorphism once the class field axiom holds; its
inverse is the arithmetic Artin map. Throughout this lesson, neither
existence of that inverse nor a reciprocity theorem is an assumption.

\subsection{2. A degree map supplies the unramified
direction}\label{a-degree-map-supplies-the-unramified-direction}

Let \(G\) be profinite and let

\[
d:G\longrightarrow\widehat{\mathbf Z}
\]

be a continuous surjection. We use field notation for its subgroup
lattice: the ground field \(k\) means \(G_k=G\), a finite extension
\(K/k\) means an open subgroup \(G_K\), and \(K\subset L\) means
\(G_L\subset G_K\). Its degree is \([L:K]=[G_K:G_L]\). Composita mean
intersections of closed subgroups. In applications these will be actual
fields and Galois groups; all the arguments here use only the indicated
subgroup operations.

Put \(I=\ker d\) and let \(\widetilde k\) denote its fixed object. For
finite \(K/k\), define

\[
f_K=[\widehat{\mathbf Z}:d(G_K)],\qquad
I_K=G_K\cap I,\qquad
d_K=f_K^{-1}d|_{G_K}.
\]

The image \(d(G_K)\) is the open subgroup \(f_K\widehat{\mathbf Z}\).
Division by \(f_K\) therefore defines a continuous surjection onto
\(\widehat{\mathbf Z}\), with kernel \(I_K\). The fixed object of
\(I_K\) is \(\widetilde K=K\widetilde k\), and

\[
\operatorname{Gal}(\widetilde K/K)\simeq\widehat{\mathbf Z}.
\]

Its element of degree 1 is the \textbf{arithmetic Frobenius}
\(\varphi_K\).

For finite \(L/K\), set

\[
f_{L/K}=[d(G_K):d(G_L)],\qquad
e_{L/K}=[I_K:I_L].
\]

Both are multiplicative in towers, and

\[
[L:K]=e_{L/K}f_{L/K},\qquad
f_L=f_Kf_{L/K},\qquad
d_K|_{G_L}=f_{L/K}d_L.
\tag{1}
\]

Indeed, \(I_KG_L\) is a subgroup of \(G_K\), and its two successive
indices are \(f_{L/K}\) and \(e_{L/K}\). This proves the first equality
without assuming that \(L/K\) is Galois. The other two follow directly
from the images of \(d\). Call \(L/K\) \textbf{unramified} when
\(e_{L/K}=1\), equivalently \(L\subset\widetilde K\). It is then cyclic
and \([L:K]=f_{L/K}\). Call it \textbf{totally ramified} when
\(f_{L/K}=1\). The largest unramified subextension of \(L/K\) is
\(L\cap\widetilde K\), of degree \(f_{L/K}\).

Fix for now a finite Galois \(L/K\). Write

\[
\Gamma=\operatorname{Gal}(\widetilde L/K),\qquad
J=\operatorname{Gal}(\widetilde L/\widetilde K).
\]

Here \(J\) is finite of order \(e_{L/K}\), and

\[
1\longrightarrow J\longrightarrow\Gamma
\xrightarrow{d_K}\widehat{\mathbf Z}\longrightarrow1.
\tag{2}
\]

Choose \(\phi\in\Gamma\) of degree 1. The map from
\(\widehat{\mathbf Z}\) to the closure of its powers, sending 1 to
\(\phi\), is an isomorphism: composing it with \(d_K\) gives the
identity. Consequently every element of \(\Gamma\) has a unique
expression \(j\phi^z\), with \(j\in J\), \(z\in\widehat{\mathbf Z}\).
This splitting is a tool for the proof, not an extra datum in the final
map.

A \textbf{positive Frobenius lift} is \(s\in\Gamma\) with \(d_K(s)=n\) a
positive ordinary integer. Its closed cyclic subgroup
\(S=\overline{\langle s\rangle}\) is a copy of \(\widehat{\mathbf Z}\):
the degree map on it is multiplication by \(n\), which is injective. In
particular \(S\cap J=1\). If \(\Sigma\) is its fixed object, then

\[
[\Sigma:K]=n|J|,\qquad f_{\Sigma/K}=n,
\qquad\widetilde\Sigma=\widetilde L,
\qquad s=\varphi_\Sigma.
\tag{3}
\]

For the degree, the cosets of \(S\) in \(\Gamma\) are represented by
\(\phi^i j\), \(0\leq i<n\), \(j\in J\). Their distinctness follows by
applying \(d_K\) and then using \(S\cap J=1\); their number is the index
given by (2). Intersecting the stabilizers with \(J\) proves the
unramified-compositum assertion. Finally
\(d_\Sigma(s)=d_K(s)/f_{\Sigma/K}=1\).

Every \(g\in\operatorname{Gal}(L/K)\) has a positive Frobenius lift.
Choose any lift to \(\Gamma\), of degree \(a\in\widehat{\mathbf Z}\).
Its residue modulo \(f_{L/K}\) is determined by \(g\)'s action on
\(L\cap\widetilde K\). Choose a positive integer \(n\) with that
residue. The kernel of \(\Gamma\to\operatorname{Gal}(L/K)\), namely
\(\operatorname{Gal}(\widetilde L/L)\), has degree image
\(f_{L/K}\widehat{\mathbf Z}\). An element of this kernel of degree
\(n-a\) adjusts the chosen lift to degree \(n\).

\subsection{3. Norms and a valuation on a
module}\label{norms-and-a-valuation-on-a-module}

Let \(A\) be a discrete abelian group with a continuous \(G\)-action. We
write it \textbf{additively} in this lesson so that the operator
calculations are readable. For a multiplicative arithmetic module, read
a sum as a product, a difference as a quotient, and zero as 1. Put
\(A_K=A^{G_K}\). For finite \(L/K\), define

\[
N_{L/K}(a)=\sum_{t\in G_K/G_L}t(a).
\tag{4}
\]

Representatives in (4) are for left cosets. Since \(a\) is fixed by
\(G_L\), their choices do not matter. Left multiplication permutes these
cosets, so the sum belongs to \(A_K\). Decomposing cosets along a tower
proves \(N_{M/K}=N_{L/K}N_{M/L}\). Conjugating the cosets proves

\[
N_{gL/gK}(ga)=gN_{L/K}(a).
\tag{5}
\]

Here \(gK\) has stabilizer \(gG_Kg^{-1}\). These identities do not
require intermediate fields to be Galois. For an infinite field object,
\(A_E\) is the union of the invariants of its finite subobjects: an
element has an open stabilizer by continuity.

Choose a subgroup \(V\subset\widehat{\mathbf Z}\) containing the
ordinary integers, such that the inclusion induces

\[
V/nV\simeq\widehat{\mathbf Z}/n\widehat{\mathbf Z}
\quad(n\geq1).
\tag{6}
\]

A \textbf{henselian valuation for the degree map} is a group
homomorphism \(v:A_k\to\widehat{\mathbf Z}\) with

\[
v(A_k)=V,\qquad v(N_{K/k}A_K)=f_KV
\quad\text{for every finite }K/k.
\tag{7}
\]

This is a condition on an arithmetic module and its norms, not a field
valuation. Allowing \(V\) to be larger than \(\mathbf Z\) matters: the
local case will use \(V=\mathbf Z\), whereas the global number-field
construction can use \(V=\widehat{\mathbf Z}\). All proofs here cover
both. Requiring only integer values would unnecessarily exclude the
latter application.

For finite \(K/k\), put

\[
v_K=f_K^{-1}vN_{K/k}:A_K\longrightarrow V.
\]

Condition (7) makes this well defined and onto. Norm transitivity and
(1) give

\[
v_K(N_{L/K}a)=f_{L/K}v_L(a).
\tag{8}
\]

For \(a\in A_K\), equation (4) gives \(N_{L/K}a=[L:K]a\); cancellation
in the torsion-free group \(\widehat{\mathbf Z}\) then yields

\[
v_L(a)=e_{L/K}v_K(a).
\tag{9}
\]

Equation (5) also gives \(v_{gK}(ga)=v_K(a)\), since a norm to \(k\) is
\(G\)-invariant. A \textbf{prime element} is any \(\pi_K\in A_K\) with
\(v_K(\pi_K)=1\); it exists by surjectivity. Let \(U_K=\ker v_K\), the
\textbf{unit module}. Units remain units under inclusion or norm. A
prime element remains prime under unramified inclusion; its norm remains
prime under a totally ramified norm.

Our only cohomological hypothesis in this lesson is the
\textbf{unramified unit axiom}:

\[
\widehat H^0(\operatorname{Gal}(E/F),U_E)=0,
\qquad
\widehat H^{-1}(\operatorname{Gal}(E/F),U_E)=0
\tag{10}
\]

for every finite unramified \(E/F\). Thus every unit downstairs is a
norm, and every norm-zero unit upstairs is a difference \((\tau-1)y\),
where \(\tau\) generates the cyclic group. The cyclic Tate definitions
in lesson 2 make these two meanings explicit. Later lessons verify (10)
in the arithmetic applications; it is an explicit hypothesis of the
present abstract theorem.

\subsection{4. Calibrating the unramified
map}\label{calibrating-the-unramified-map}

Before treating ramification, check what the general hypotheses say
about an unramified extension. If its degree is \(n\), the valuation of
its norm image is \(nV\), and the unit axiom removes the remaining
kernel. Thus a prime class has to be the Frobenius class. The following
proposition gives the precise statement, including the value group
\(V=\widehat{\mathbf Z}\) needed later for number fields.

\subsubsection{Proposition 4.3. Unramified
reciprocity}\label{proposition-4.3.-unramified-reciprocity}

If \(L/K\) is unramified of degree \(n\), define \(r_{L/K}\) by sending
its arithmetic Frobenius generator to a prime class. This is independent
of the prime, and

\[
r_{L/K}(\varphi_{L/K})=\pi_K\pmod{N_{L/K}A_L},
\]

and \(r_{L/K}\) is an isomorphism.

\textbf{Proof.} Two primes differ by a unit, which is a norm by (10), so
they define the same class. Its order divides \(n\), since
\(n\pi_K=N_{L/K}(\pi_K)\); hence the generator prescription defines a
homomorphism.

The valuation induces a surjection \(A_K/N_{L/K}A_L\to V/nV\), by (8).
It is injective as well. If \(v_K(a)=nb\), choose \(c\in A_L\) with
\(v_L(c)=b\). Then \(a-N_{L/K}c\) is a unit in \(U_K\), hence a norm
from \(U_L\) by (10); thus \(a\) is a norm. By (6),
\(V/nV\simeq\mathbf Z/n\mathbf Z\), and the image of \(\pi_K\) is its
generator 1. The cyclic Galois group is also generated by
\(\varphi_{L/K}\) and has order \(n\). The map between them sending
generator to generator is an isomorphism. \(\square\)

The norm-of-prime construction in section 5 specializes to this map:
here \(\widetilde L=\widetilde K\), and the degree-one lift
\(\varphi_K\) has fixed field \(K\), so (12) gives \(\pi_K\).

This argument does not assume that an arbitrary valuation is an ordinary
integer, or that an element has an expression \(u+m\pi_K\) with
\(m\in\mathbf Z\). Such an expression need not exist when
\(V=\widehat{\mathbf Z}\). Lifting the divisible valuation by a norm is
the step that treats both value groups uniformly.

\subsection{5. Building the candidate
class}\label{building-the-candidate-class}

Choose a positive lift of the desired finite automorphism. Its fixed
field is finite by the degree calculation in section 2. The norm
identity below tells us how to compare candidate classes without
assuming that their prime elements belong to the same field.

Every finite subextension of \(\widetilde L/L\) is unramified. Equation
(9) therefore defines a consistent valuation \(w\) on
\(A_{\widetilde L}\), by restricting to a finite subextension containing
\(L\). Its kernel is \(U_{\widetilde L}\). The action of \(\Gamma\)
preserves \(w\), by (5). Each prime \(\pi_\Sigma\) from (3) has
\(w(\pi_\Sigma)=1\): both \(\Sigma\) and \(L\) have ramification index
\(|J|\) over \(K\), so their compositum is unramified over each.

Let \(N_J=\sum_{j\in J}j\), and put
\(P_n(\phi)=1+\phi+\cdots+\phi^{n-1}\). Normality of \(J\) implies that
\(N_J\) commutes with \(\phi\). The coset representatives used in (3)
prove the \textbf{norm identity}

\[
N_{\Sigma/K}(a)=N_JP_n(\phi)a
\quad(a\in A_\Sigma, d_K(s)=n).
\tag{11}
\]

For a positive lift \(s\), let \(\Sigma\) be its fixed field and choose
a prime \(\pi_\Sigma\). Define provisionally

\[
R(s)=N_{\Sigma/K}(\pi_\Sigma)\pmod{N_{L/K}A_L}.
\tag{12}
\]

Changing the prime by \(u\in U_\Sigma\) does not change (12). Indeed,
\(L\Sigma/\Sigma\) is unramified, so (10) gives
\(u=N_{L\Sigma/\Sigma}b\) for a unit \(b\). Thus
\(N_{\Sigma/K}u=N_{L/K}N_{L\Sigma/L}b\) is a norm from \(L\). This
already handles the first choice.

\subsection{6. The unit defect and
composition}\label{the-unit-defect-and-composition}

The unresolved issue is now a calculation about units. Two prime
elements in the enlarged tower have equal valuation, so their difference
in additive notation is a unit. For a product of lifts, the valuation
contributions also cancel. What must be proved is that this residual
unit becomes a norm after averaging over inertia. The next lemma is
designed for exactly that defect; both halves of the unit axiom are used
in its proof.

We now prove the unit statement that controls composition. Write
\(I_JU_{\widetilde L}\) for the subgroup generated by \((j-1)u\), with
\(j\in J\), \(u\in U_{\widetilde L}\).

\textbf{Unit descent lemma.} If \(u\in U_{\widetilde L}\) and
\((\phi-1)u\in I_JU_{\widetilde L}\), then

\[
N_Ju\in N_{M/K}U_M
\]

for every finite \(M/K\) contained in \(\widetilde L\).

\textbf{Proof.} Express \((\phi-1)u=\sum_i(j_i-1)u_i\), using finitely
many units. Enlarge the given \(M\) to a finite Galois extension of
\(K\) inside \(\widetilde L\) containing \(L\), \(u\), and all \(u_i\).
Norm transitivity will imply the conclusion for the original \(M\), so
retain the enlarged name. Put \(m=[M:K]\) and \(\psi=\phi^m\). This
element fixes \(M\). It commutes with \(J\): conjugation by \(\psi\) is
trivial on the restrictions to \(L\), and restriction embeds \(J\) in
\(\operatorname{Gal}(L/K)\). It also commutes with \(\phi\), so it is
central in \(\Gamma\).

The closure of the powers of \(\phi\) maps isomorphically to
\(\widehat{\mathbf Z}\) under \(d_K\). To justify this, the map
\(a\mapsto\phi^a\) is defined in every finite quotient by compatible
integer powers, and their inverse limit defines it on
\(\widehat{\mathbf Z}\). Its composite with \(d_K\) is the identity. Its
image is exactly that closure, since ordinary integers are dense. Thus
it intersects \(J\) trivially. Every element of \(\Gamma\) is a product
\(j\phi^a\), by first matching its degree. In particular the closed
subgroup generated by \(\psi=\phi^m\) has index \(m|J|\), and the
subgroup generated by \(\psi^m\) has index \(m^2|J|\).

Let \(E\) and \(E_m\) be their fixed fields. Centrality makes both
closed subgroups normal, so these are finite Galois extensions of \(K\).
They contain \(M\), since \(\psi\) fixes \(M\). Their relative group is
\(m\widehat{\mathbf Z}/m^2\widehat{\mathbf Z}\), generated by \(\psi\),
and its intersection with inertia is trivial. Hence \(E_m/E\) is
unramified cyclic of degree \(m\). Axiom (10) gives units
\(b,b_i\in U_{E_m}\) with

\[
P_m(\psi)b=u,\qquad P_m(\psi)b_i=u_i.
\]

The norm of \((\phi-1)b-\sum_i(j_i-1)b_i\) is zero, since \(\psi\)
commutes with \(\phi,J\). The second half of (10) therefore supplies
\(y\in U_{E_m}\) such that

\[
(\phi-1)b-\sum_i(j_i-1)b_i
=(\psi-1)y=(\phi-1)P_m(\phi)y.
\]

Apply \(N_J\). It kills the \(J\)-differences, giving

\[
(\phi-1)\bigl(N_Jb-P_m(\phi)N_Jy\bigr)=0.
\]

The parenthesized unit is fixed by \(J\) and by \(\phi\), hence by all
of \(\Gamma\); call it \(z\in U_K\). Apply \(P_m(\psi)\) to
\(N_Jb=P_m(\phi)N_Jy+z\), and put \(y'=P_m(\psi)y\in U_E\). We obtain

\[
N_Ju=N_JP_m(\phi)y'+mz
=N_{E/K}(y')+N_{M/K}(z).
\]

The first equality in the last step uses (11), and the second uses
\(z\in U_K\). Since \(M\subset E\), both terms are norms from \(U_M\).
This proves the assertion. \(\square\)

Both halves of the unit axiom were used: one lifts the units by norms,
and the other lifts the relation by a difference. Surjectivity of unit
norms alone would not prove this lemma.

\textbf{Composition lemma.} For positive lifts \(s_1,s_2\),

\[
R(s_1s_2)=R(s_1)+R(s_2).
\tag{13}
\]

\textbf{Proof.} Write \(s_3=s_1s_2\), \(n_i=d_K(s_i)\), so
\(n_3=n_1+n_2\), and choose corresponding primes \(\pi_i\). Set

\[
s_4=\phi^{n_1}s_2\phi^{-n_1},\qquad
\pi_4=\phi^{n_1}\pi_2,\qquad n_4=n_2.
\]

The norm-conjugation identity (5) gives \(R(s_4)=R(s_2)\), even before
taking classes, since the resulting norm belongs to \(A_K\). Put

\[
t_i=\phi^{n_i}s_i^{-1}\in J.
\]

A direct multiplication gives \(t_3=t_4t_1\). Also \(s_i\pi_i=\pi_i\),
so \((\phi^{n_i}-1)\pi_i=(t_i-1)\pi_i\). Consider

\[
u=P_{n_3}(\phi)\pi_3-P_{n_1}(\phi)\pi_1
-P_{n_2}(\phi)\pi_4.
\]

It is a unit because its valuation is \(n_3-n_1-n_2=0\). Let
\(a=\pi_3-\pi_4\) and \(b=\pi_1-\pi_4\), also units. The identity
\((\phi-1)P_n(\phi)=\phi^n-1\) gives

\[
\begin{aligned}
(\phi-1)u
&=(t_3-1)\pi_3-(t_1-1)\pi_1-(t_4-1)\pi_4\\
&=(t_3-1)a-(t_1-1)b+(t_4-1)(t_1-1)\pi_4.
\end{aligned}
\tag{14}
\]

The element \((t_1-1)\pi_4\) is a unit since \(t_1\) preserves
valuation. Every term on the last line of (14) thus belongs to
\(I_JU_{\widetilde L}\). The unit descent lemma makes \(N_Ju\) a norm
from \(U_L\). By (11), its class is exactly \(R(s_3)-R(s_1)-R(s_4)\),
proving (13). \(\square\)

\subsubsection{Proposition 4.1. Independence of the
choices}\label{proposition-4.1.-independence-of-the-choices}

The value (12) depends only on \(s|_L\). It is independent of the
positive lift and of its prime element.

\textbf{Proof.} Prime independence was proved just before the
composition lemma. Suppose \(s,s'\) have the same restriction to \(L\).
If their degrees are equal, they are equal: the actions on \(L\) and
\(\widetilde K\) determine an action on their compositum
\(\widetilde L\). If their degrees differ, interchange them if necessary
so that \(d_K(s')>d_K(s)\), and set \(h=s^{-1}s'\). This is a positive
lift restricting to the identity on \(L\). Its fixed field contains
\(L\), so its defining norm belongs to \(N_{L/K}A_L\), and \(R(h)=0\).
Equation (13) now gives \(R(s')=R(s)+R(h)=R(s)\). \(\square\)

\subsubsection{Proposition 4.2. The reciprocity
homomorphism}\label{proposition-4.2.-the-reciprocity-homomorphism}

The prescription

\[
r_{L/K}(g)=R(s),\qquad s|_L=g,
\]

defines a homomorphism

\[
r_{L/K}:\operatorname{Gal}(L/K)^{\mathrm{ab}}
\longrightarrow A_K/N_{L/K}A_L.
\tag{15}
\]

For a multiplicatively written \(A\), equation (13) says exactly that
\(r\) is multiplicative.

\textbf{Proof.} Positive lifts exist by section 2, and Proposition 4.1
makes the prescription independent of them. The product of lifts of
\(g_1,g_2\) is a positive lift of \(g_1g_2\), so (13) proves the
homomorphism property. The target is abelian, so the homomorphism kills
commutators and factors through the finite group's abelianization.
\(\square\)

\subsection{7. Norms, conjugation and
transfer}\label{norms-conjugation-and-transfer}

\subsubsection{Proposition 4.4.
Functoriality}\label{proposition-4.4.-functoriality}

The maps (15) have the following three compatibilities.

\begin{enumerate}
\def\labelenumi{\arabic{enumi}.}
\tightlist
\item
  If \(K\subset K'\), \(L\subset L'\), and both \(L/K\), \(L'/K'\) are
  finite Galois, then \[
  N_{K'/K}\,r_{L'/K'}(g')=r_{L/K}(g'|_L).
  \tag{16}
  \] The norm on the left is taken on the corresponding quotient
  modules.
\item
  For \(g\in G\), \[
  r_{gL/gK}(g\sigma g^{-1})=g\,r_{L/K}(\sigma).
  \tag{17}
  \]
\item
  If \(K'\) is intermediate in a finite Galois \(L/K\), the inclusion of
  modules satisfies \[
  r_{L/K}(\sigma)\text{ viewed in }A_{K'}/N_{L/K'}A_L
  =r_{L/K'}(\operatorname{Ver}\sigma).
  \tag{18}
  \] Here \(\operatorname{Ver}:\operatorname{Gal}(L/K)^{\mathrm{ab}}\to
  \operatorname{Gal}(L/K')^{\mathrm{ab}}\) is transfer.
\end{enumerate}

\textbf{Proof of (16).} Lift \(g'\) to a positive
\(s'\in\operatorname{Gal}(\widetilde L'/K')\), of degree \(n'\). Its
restriction \(s\) to \(\widetilde L\) is positive, of degree
\(f_{K'/K}n'\). If \(\Sigma'\) and \(\Sigma\) are their fixed fields,
then \(\Sigma=\Sigma'\cap\widetilde L\). Their absolute residue degrees
are both \(f_{K'}n'\), so \(f_{\Sigma'/\Sigma}=1\). For a prime
\(\pi'\in A_{\Sigma'}\), equation (8) makes
\(\pi=N_{\Sigma'/\Sigma}\pi'\) a prime of \(\Sigma\). Norm transitivity
gives the exact equality

\[
N_{\Sigma/K}\pi
=N_{\Sigma'/K}\pi'
=N_{K'/K}N_{\Sigma'/K'}\pi'.
\]

This gives (16). The quotient norm is well defined because
\(N_{K'/K}N_{L'/K'}A_{L'}=N_{L'/K}A_{L'}\subset N_{L/K}A_L\).

\textbf{Proof of (17).} A lift \(s\) with fixed field \(\Sigma\) becomes
\(gsg^{-1}\) with fixed field \(g\Sigma\). Its degree remains the same,
since the degree map has abelian target. The element \(g\pi_\Sigma\) is
prime by the conjugation identity for valuations. Equation (5) gives
(17) directly.

For (18), we first establish the needed group-theoretic formula. Let
\(T\) be finite and \(H\subset T\). Choose representatives \(R\) for
left cosets \(T/H\). For \(\sigma\in T\), write

\[
\sigma r=r' h_r(\sigma),\qquad r'\in R,\ h_r(\sigma)\in H.
\]

Then \(\operatorname{Ver}(\sigma)=\prod_{r\in R}h_r(\sigma)\) in
\(H^{\mathrm{ab}}\). Replacing \(r\) by \(rk_r\) replaces its factor by
\(k_{r'}^{-1}h_r(\sigma)k_r\); the \(k_r\) cancel in the abelian
product. Thus transfer does not depend on representatives. The factors
for \(\sigma\tau\) are those for \(\sigma\) at the cosets permuted by
\(\tau\), followed by those for \(\tau\). Their products prove that
transfer is a homomorphism, hence factors through \(T^{\mathrm{ab}}\).
On an orbit of \(\sigma\) in \(T/H\), choose representatives
\(r,\sigma r,\ldots,\sigma^{a-1}r\). The contribution is
\(r^{-1}\sigma^a r\), where \(a\) is that orbit's length. Therefore

\[
\operatorname{Ver}(\sigma)=\prod_r r^{-1}\sigma^{a_r}r
\quad\text{in }H^{\mathrm{ab}},
\tag{19}
\]

with one representative per orbit.

Now take \(T=\operatorname{Gal}(L/K)\), \(H=\operatorname{Gal}(L/K')\).
The inclusion in (18) is well defined: decomposing \(T\) into right
cosets \(H t\) shows that \(N_{L/K}a=N_{L/K'}(\sum_t t a)\), so every
norm from \(L/K\) is a norm from \(L/K'\).

Choose a positive lift \(s\in\Gamma\) of \(\sigma\), with fixed field
\(\Sigma\), and write \(S=\overline{\langle s\rangle}\),
\(\Gamma'=\operatorname{Gal}(\widetilde L/K')\). The finite coset spaces
\(\Gamma/\Gamma'\) and \(T/H\) agree. Representatives for the
\(S\)-orbits may therefore be chosen to lift those in (19). For each,
put

\[
s_r=r^{-1}s^{a_r}r,\qquad
S_r=\Gamma'\cap r^{-1}Sr=\overline{\langle s_r\rangle}.
\]

The equality follows because the open subgroup of the procyclic group
\(S\) stabilizing that coset has index \(a_r\). The element \(s_r\) is a
positive lift over \(K'\): its degree is \(a_r d_K(s)/f_{K'/K}\), a
positive integer since its restriction belongs to \(\Gamma'\). Let
\(\Sigma_r\) be its fixed field. Then

\[
\Sigma_r=(r^{-1}\Sigma)K',
\qquad \Sigma_r/(r^{-1}\Sigma)\text{ is unramified of degree }a_r.
\]

Thus \(\pi_r=r^{-1}\pi_\Sigma\) is still prime in \(\Sigma_r\). Finally,
the double cosets \(\Gamma'\backslash\Gamma/S\) partition the cosets
used in the norm from \(\Sigma\). The part represented by \(r^{-1}\) has
stabilizer \(S_r\) in \(\Gamma'\), so summing it gives

\[
N_{\Sigma/K}\pi_\Sigma
=\sum_r N_{\Sigma_r/K'}\pi_r
\quad\text{in }A_{K'}.
\tag{20}
\]

Each summand represents \(r_{L/K'}(r^{-1}\sigma^{a_r}r)\). Since
reciprocity is a homomorphism, (19) and (20) prove (18). \(\square\)

The arrows have different arithmetic meanings. Restricting a Galois
automorphism corresponds to taking a norm. Transfer corresponds to
including an arithmetic element in a larger field. Confusing those two
operations reverses a reciprocity diagram.

\subsection{8. Two complete degree
models}\label{two-complete-degree-models}

The simplest complete model isolates the degree contribution. After
checking it, compare it with the local-field application: the degree is
still present, but units now carry additional information and the
descent lemma controls that information.

Take
\(G=\operatorname{Gal}(\overline{\mathbf F}_q/\mathbf F_q)=\widehat{\mathbf Z}\),
with \(d\) the identity. Let \(A=\mathbf Z\) have trivial action, and
let \(v\) be the inclusion in \(\widehat{\mathbf Z}\). Every \(A_K\) is
\(\mathbf Z\). For \(K=\mathbf F_{q^m}\), the norm to the ground field
is multiplication by \(m\), so (7) holds with \(V=\mathbf Z\), and
\(v_K\) is the identity. Every extension is unramified and every unit
module is zero, so (10) holds.

For \(L/K\) of degree \(n\), the quotient module is
\(\mathbf Z/n\mathbf Z\), and reciprocity sends arithmetic Frobenius to
1. A positive lift of degree \(b\) has fixed field of degree \(b\) over
\(K\); the norm of its prime element 1 is \(b\), as required. This is
the degree formation associated with finite fields. It uses
\(\mathbf Z\), not the multiplicative group of the finite field, whose
norm is surjective and could not give this quotient.

For local fields the arithmetic module will instead be the
multiplicative group of a separable closure, and the degree direction
will be the unramified tower. The local reciprocity lesson will verify
the valuation and unit hypotheses. The construction above already
identifies the map once those hypotheses are verified.

The same model can be built with \(A=\widehat{\mathbf Z}\), still
treated as a discrete module with trivial action, and \(v\) the
identity. A degree-\(m\) norm is multiplication by \(m\), so its image
is \(m\widehat{\mathbf Z}\). All normalized valuations are the identity,
the unit modules are zero, and the quotient for degree \(n\) is again
\(\mathbf Z/n\mathbf Z\). An element of \(\widehat{\mathbf Z}\) need not
be an ordinary integral multiple of 1, but its class in this finite
quotient is. This explains why the finite-quotient condition on \(V\),
rather than an integer-valued hypothesis, is the correct input to the
proof.

\subsection{Exercises}\label{exercises}

\begin{enumerate}
\def\labelenumi{\arabic{enumi}.}
\tightlist
\item
  \textbf{Easy.} Work out the degree formation of section 8, including
  its normalized valuations, norm groups and reciprocity for every
  \(\mathbf F_{q^{mn}}/\mathbf F_{q^m}\).
\item
  \textbf{Medium.} Prove that changing the prime of a positive lift's
  fixed field does not change its reciprocity class. Identify the
  unramified extension on which the unit axiom is applied.
\item
  \textbf{Medium.} Prove conjugation compatibility (17) directly from
  the norm definition, including preservation of prime elements.
\item
  \textbf{Hard.} Starting with the unit descent lemma, prove the
  composition formula (13). Check explicitly the identity \(t_3=t_4t_1\)
  and why every term of (14) is a difference of units.
\end{enumerate}

\subsection{Solutions}\label{solutions}

\begin{enumerate}
\def\labelenumi{\arabic{enumi}.}
\tightlist
\item
  The subgroup for \(\mathbf F_{q^m}\) is \(m\widehat{\mathbf Z}\), so
  \(f_K=m\), \(d_K=d/m\), and \(I_K=0\). Trivial action gives
  \(A_K=\mathbf Z\). A norm along a degree-\(n\) extension is
  multiplication by \(n\), hence its group is \(n\mathbf Z\). Since
  \(N_{K/k}\) is multiplication by \(m\), \(v_K=m^{-1}vN_{K/k}\) is the
  identity. The only prime element is 1, the unit module is zero, and
  both Tate groups in the unit axiom vanish. The cyclic Galois group of
  order \(n\) maps isomorphically to \(\mathbf Z/n\mathbf Z\), with
  Frobenius going to 1 and its \(b\)-th power to \(b\bmod n\).
\item
  Two primes differ by \(u\in U_\Sigma\). Equation (3) puts \(L\Sigma\)
  inside \(\widetilde\Sigma\), so \(L\Sigma/\Sigma\) is finite
  unramified cyclic. The norm-surjectivity half of (10) gives
  \(u=N_{L\Sigma/\Sigma}b\), with \(b\) a unit. Taking its norm to \(K\)
  gives \(N_{\Sigma/K}u=N_{L/K}N_{L\Sigma/L}b\), which is zero in the
  quotient by norms from \(L\). The extension required is
  \(L\Sigma/\Sigma\), not a possibly ramified extension \(L/K\).
\item
  Representatives \(t\) for \(G_K/G_\Sigma\) become \(gtg^{-1}\) for
  \(G_{gK}/G_{g\Sigma}\), so the norm of \(ga\) is \(gN_{\Sigma/K}a\).
  The same identity for the norm to \(k\), together with unchanged
  \(f_\Sigma\), yields \(v_{g\Sigma}(ga)=v_\Sigma(a)\). Thus
  \(g\pi_\Sigma\) is prime. A positive lift becomes \(gsg^{-1}\), of the
  same degree and with fixed field \(g\Sigma\). Applying its defining
  norm proves (17), and conjugation also carries the denominator norm
  subgroup to the corresponding denominator.
\item
  Choose \(s_4=\phi^{n_1}s_2\phi^{-n_1}\), \(\pi_4=\phi^{n_1}\pi_2\).
  Then \[
  t_4t_1
  =\phi^{n_2}\phi^{n_1}s_2^{-1}\phi^{-n_1}
    \phi^{n_1}s_1^{-1}
  =\phi^{n_1+n_2}(s_1s_2)^{-1}=t_3.
  \] Form \(u\) as in section 6. Each prime has valuation 1, so
  \(w(u)=0\). Applying \(\phi-1\) replaces \(P_{n_i}(\phi)\pi_i\) by
  \((t_i-1)\pi_i\). Substitute \(\pi_3=a+\pi_4\), \(\pi_1=b+\pi_4\). The
  remaining coefficient of \(\pi_4\) is
  \(t_3-t_1-t_4+1=(t_4-1)(t_1-1)\), giving (14). Here \(a,b\), and
  \((t_1-1)\pi_4\) all have valuation zero. Thus
  \((\phi-1)u\in I_JU_{\widetilde L}\). Unit descent makes \(N_Ju\) a
  norm from \(L\); (11) identifies its class with
  \(R(s_1s_2)-R(s_1)-R(s_4)\), and (5) identifies \(R(s_4)=R(s_2)\).
  This proves composition with the chosen left-action convention.
\end{enumerate}

\subsection{Editable edition}\label{editable-edition}

The reading edition provides the complete LaTeX source of this lesson,
the cumulative course LaTeX and the editable source ZIP. The archive
contains all twenty-four Markdown lessons, complete LaTeX bodies,
original diagrams, metadata and reproduction instructions.

\subsection{References}\label{references}

Sections 2--3 give the full abstract degree and valuation hypotheses,
including the general permitted value group. Sections 5--7 prove unit
descent, independence, composition and all three functorialities. The
central finite unramified tower used in descent is constructed
explicitly.

\begin{itemize}
\tightlist
\item
  \href{https://kskedlaya.org/cft/sec_abstractcft1.html}{Kiran S.
  Kedlaya, Notes on class field theory, author-hosted HTML edition}.
\end{itemize}

The abstract hypotheses allow the full value group specified in section
3. The proof does not assume a splitting of that valuation sequence.

The \href{../free-proofs.html}{proof guide} gives the lesson sequence
and the exact prerequisite record. External references accompany the
written arguments.

\end{document}
